% =============================================================================
% Common Preamble for Network+ Lecture Series
% This file provides consistent formatting, colors, packages, and styles
% across all lecture modules for both beamer presentations and article mode
% =============================================================================

% =============================================================================
% THEME AND COLORS
% =============================================================================
\usetheme{Madrid}
\usecolortheme{whale}

% Custom colors - standardized across all lectures
\definecolor{rctcblue}{RGB}{0, 74, 136}
\definecolor{rctcgold}{RGB}{255, 199, 44}
\definecolor{networkgreen}{RGB}{0, 128, 64}
\definecolor{alertred}{RGB}{180, 0, 0}
\definecolor{aborange}{RGB}{230, 126, 34}
\definecolor{copperbrown}{RGB}{184, 115, 51}
\definecolor{faborange}{RGB}{255, 165, 0}
\definecolor{fiberyellow}{RGB}{255, 215, 0}
\definecolor{shirepurple}{RGB}{102, 51, 153}
\definecolor{ipblue}{RGB}{52, 152, 219}
\definecolor{ippurple}{RGB}{155, 89, 182}
\definecolor{vlangreen}{RGB}{39, 174, 96}
\definecolor{tcpblue}{RGB}{30, 100, 180}
\definecolor{udpgreen}{RGB}{50, 160, 80}
\definecolor{dhcppurple}{RGB}{120, 60, 160}
\definecolor{dnsyellow}{RGB}{200, 170, 50}
\definecolor{batblack}{RGB}{30, 30, 40}
\definecolor{gothamgray}{RGB}{80, 85, 95}
\definecolor{paddysgreen}{RGB}{19, 77, 33}
\definecolor{beerlight}{RGB}{255, 199, 44}
\definecolor{sunnyorange}{RGB}{255, 140, 0}
\definecolor{ogregreen}{RGB}{34, 139, 34}
\definecolor{swampbrown}{RGB}{101, 67, 33}

% Set theme colors
\setbeamercolor{palette primary}{bg=rctcblue,fg=white}
\setbeamercolor{palette secondary}{bg=rctcblue!80,fg=white}
\setbeamercolor{palette tertiary}{bg=rctcblue!60,fg=white}
\setbeamercolor{structure}{fg=rctcblue}
\setbeamercolor{title}{fg=white,bg=rctcblue}
\setbeamercolor{frametitle}{fg=white,bg=rctcblue}
\setbeamercolor{block title}{bg=rctcblue,fg=white}
\setbeamercolor{block body}{bg=rctcblue!10,fg=black}
\setbeamercolor{block title alerted}{bg=alertred,fg=white}
\setbeamercolor{block body alerted}{bg=alertred!10,fg=black}
\setbeamercolor{block title example}{bg=networkgreen,fg=white}
\setbeamercolor{block body example}{bg=networkgreen!10,fg=black}

% =============================================================================
% PACKAGES
% =============================================================================
\usepackage[utf8]{inputenc}
\usepackage[T1]{fontenc}
\usepackage{lmodern}  % Better font rendering
\usepackage{graphicx}
\usepackage{tikz}
\usepackage{booktabs}
\usepackage{array}
\usepackage{multirow}
\usepackage{hyperref}
\usepackage{adjustbox}
\usepackage{amssymb}
\usepackage{amsmath}
\usepackage{calc}
\usepackage{fontawesome5}
% xcolor with [table] option: avoid colortbl.4ht hook conflict in tex4ht
\ifdefined\HCode
  \usepackage{xcolor}
  \usepackage{colortbl}  % provides \rowcolor etc.
\else
  \usepackage[table]{xcolor}  % includes colortbl for pdflatex
\fi

% =============================================================================
% TIKZ LIBRARIES AND SETUP
% =============================================================================
\usetikzlibrary{
    shapes.geometric,
    shapes.symbols,
    shapes.misc,
    arrows.meta,
    positioning,
    calc,
    fit,
    backgrounds,
    decorations.pathreplacing,
    decorations.pathmorphing,
    matrix,
    chains
}
% shadows library causes pgfuseshading errors with tex4ht;
% only load it when NOT building HTML
\ifdefined\HCode\else
    \usetikzlibrary{shadows}
\fi

% Graphics path for icons and chapter images
\graphicspath{
    {../images/icons/}
    {./images/icons/}
    {images/icons/}
    {../images/ch1/}
    {./images/ch1/}
    {images/ch1/}
    {../images/ch2/}
    {./images/ch2/}
    {images/ch2/}
    {../images/ch3/}
    {./images/ch3/}
    {images/ch3/}
    {../images/ch4/}
    {./images/ch4/}
    {images/ch4/}
    {../images/ch5/}
    {./images/ch5/}
    {images/ch5/}
    {../images/ch6/}
    {./images/ch6/}
    {images/ch6/}
    {../images/ch7/}
    {./images/ch7/}
    {images/ch7/}
    {../images/ch8/}
    {./images/ch8/}
    {images/ch8/}
    {../images/ch9/}
    {./images/ch9/}
    {images/ch9/}
    {../images/}
    {./images/}
    {images/}
}

% =============================================================================
% NETWORK DIAGRAM STYLES
% =============================================================================
\tikzset{
    % Node styles for network devices
    device/.style={
        inner sep=2pt,
        label distance=2pt
    },
    pc/.style={device},
    laptop/.style={device},
    server/.style={device},
    router/.style={device},
    switch/.style={device},
    firewall/.style={device},
    cloud/.style={device},
    phone/.style={device},
    tablet/.style={device},
    wifi/.style={device},
    % Connection styles
    ethernet/.style={
        draw=rctcblue,
        line width=1.5pt,
        -{Stealth[length=3mm]}
    },
    ethernet noarrow/.style={
        draw=rctcblue,
        line width=1.5pt
    },
    wireless/.style={
        draw=aborange,
        line width=1.5pt,
        dashed,
        -{Stealth[length=3mm]}
    },
    wireless noarrow/.style={
        draw=aborange,
        line width=1.5pt,
        dashed
    },
    wan/.style={
        draw=alertred,
        line width=2pt,
        -{Stealth[length=3mm]}
    },
    wan noarrow/.style={
        draw=alertred,
        line width=2pt
    },
    copper/.style={
        draw=copperbrown,
        line width=2pt,
        -{Stealth[length=3mm]}
    },
    copper noarrow/.style={
        draw=copperbrown,
        line width=2pt
    },
    fiber/.style={
        draw=fiberyellow,
        line width=2pt,
        -{Stealth[length=3mm]}
    },
    fiber noarrow/.style={
        draw=fiberyellow,
        line width=2pt
    },
    trunk/.style={
        draw=shirepurple,
        line width=2pt,
        -{Stealth[length=3mm]}
    },
    trunk noarrow/.style={
        draw=shirepurple,
        line width=2pt
    },
    % Packet/segment styles
    pktseg/.style={
        draw=rctcblue,
        fill=rctcblue!10,
        minimum height=0.5cm,
        font=\tiny,
        rounded corners=2pt
    },
    pktseg tcp/.style={
        draw=tcpblue,
        fill=tcpblue!15,
        minimum height=0.5cm,
        font=\tiny,
        rounded corners=2pt
    },
    pktseg udp/.style={
        draw=udpgreen,
        fill=udpgreen!15,
        minimum height=0.5cm,
        font=\tiny,
        rounded corners=2pt
    },
    pktseg highlight/.style={
        draw=aborange,
        fill=aborange!20,
        minimum height=0.5cm,
        font=\tiny,
        rounded corners=2pt
    },
    % Box styles
    ipbox/.style={
        draw=rctcblue,
        fill=rctcblue!10,
        rounded corners,
        minimum height=0.4cm,
        font=\small
    },
    portbox/.style={
        draw=tcpblue,
        fill=tcpblue!10,
        rounded corners,
        minimum height=0.4cm,
        font=\small
    },
    dnsbox/.style={
        draw=dnsyellow,
        fill=dnsyellow!15,
        rounded corners,
        minimum height=0.4cm,
        font=\small
    },
    routetable/.style={
        draw=rctcblue,
        fill=rctcblue!10,
        rounded corners,
        minimum height=0.4cm,
        font=\small
    },
    macbox/.style={
        draw=networkgreen,
        fill=networkgreen!10,
        rounded corners,
        minimum height=0.4cm,
        font=\small
    },
    % Frame segment styles (Ethernet frame diagrams)
    frameseg/.style={
        draw=rctcblue,
        fill=rctcblue!10,
        minimum height=0.5cm,
        font=\tiny,
        rounded corners=2pt
    },
    frameseg highlight/.style={
        draw=aborange,
        fill=aborange!20,
        minimum height=0.5cm,
        font=\tiny,
        rounded corners=2pt
    },
    % Process/flow styles
    process/.style={
        draw=rctcblue,
        fill=rctcblue!15,
        rounded corners,
        minimum width=2cm,
        minimum height=0.6cm,
        font=\small,
        align=center
    },
    decision/.style={
        draw=aborange,
        fill=aborange!15,
        diamond,
        aspect=2,
        minimum width=1.5cm,
        font=\small,
        align=center
    },
    % OSI Layer styles
    osilayer/.style={
        draw=rctcblue,
        fill=rctcblue!10,
        minimum width=3cm,
        minimum height=0.45cm,
        font=\small
    },
    osilayer highlight/.style={
        draw=aborange,
        fill=aborange!20,
        line width=1.5pt,
        minimum width=3cm,
        minimum height=0.45cm,
        font=\small\bfseries
    }
}

% =============================================================================
% ICON COMMANDS - Standardized using PNG files and FontAwesome
% =============================================================================
% Network device icons (PNG)
\newcommand{\iconpc}[1][0.6cm]{\resizebox{#1}{!}{\includegraphics{icon_pc}}}
\newcommand{\iconlaptop}[1][0.6cm]{\resizebox{#1}{!}{\includegraphics{icon_laptop}}}
\newcommand{\iconserver}[1][0.6cm]{\resizebox{#1}{!}{\includegraphics{icon_server}}}
\newcommand{\iconrouter}[1][0.6cm]{\resizebox{#1}{!}{
\includegraphics{icon_router}}}
\newcommand{\iconswitch}[1][0.6cm]{\resizebox{#1}{!}{
\includegraphics{icon_switch}}}
\newcommand{\iconfirewall}[1][0.6cm]{\resizebox{#1}{!}{
\includegraphics{icon_firewall}}}
\newcommand{\iconcloud}[1][0.6cm]{\resizebox{#1}{!}{
\includegraphics{icon_cloud}}}
\newcommand{\iconphone}[1][0.6cm]{\resizebox{#1}{!}{\includegraphics{icon_phone}}}
\newcommand{\icontablet}[1][0.6cm]{\resizebox{#1}{!}{\includegraphics{icon_tablet}}}
\newcommand{\iconwifi}[1][0.6cm]{\resizebox{#1}{!}{
\includegraphics{icon_wifi}}}
\newcommand{\icondatabase}[1][0.6cm]{\resizebox{#1}{!}{
\includegraphics{icon_db}}}
\newcommand{\iconlock}[1][0.6cm]{\resizebox{#1}{!}{
\includegraphics{icon_lock}}}

% For HTML/tex4ht builds: replace PNG icons with FontAwesome5 glyphs.
% dvisvgm embeds font glyphs as vectors but silently ignores em:graph (PNG)
% specials, so PNG icons would be blank in the output SVGs. FA5 glyphs work.
\ifdefined\HCode
  \renewcommand{\iconpc}[1][0.6cm]{\resizebox{#1}{!}{\faIcon{desktop}}}
  \renewcommand{\iconlaptop}[1][0.6cm]{\resizebox{#1}{!}{\faIcon{laptop}}}
  \renewcommand{\iconserver}[1][0.6cm]{\resizebox{#1}{!}{\faIcon{server}}}
  \renewcommand{\iconrouter}[1][0.6cm]{\resizebox{#1}{!}{\faIcon{route}}}
  \renewcommand{\iconswitch}[1][0.6cm]{\resizebox{#1}{!}{\faIcon{network-wired}}}
  \renewcommand{\iconfirewall}[1][0.6cm]{\resizebox{#1}{!}{\faIcon{shield-alt}}}
  \renewcommand{\iconcloud}[1][0.6cm]{\resizebox{#1}{!}{\faIcon{cloud}}}
  \renewcommand{\iconphone}[1][0.6cm]{\resizebox{#1}{!}{\faIcon{mobile-alt}}}
  \renewcommand{\icontablet}[1][0.6cm]{\resizebox{#1}{!}{\faIcon{tablet-alt}}}
  \renewcommand{\iconwifi}[1][0.6cm]{\resizebox{#1}{!}{\faIcon{wifi}}}
  \renewcommand{\icondatabase}[1][0.6cm]{\resizebox{#1}{!}{\faIcon{database}}}
  \renewcommand{\iconlock}[1][0.6cm]{\resizebox{#1}{!}{\faIcon{lock}}}
\fi

% FontAwesome icons
\newcommand{\iconuser}[1][0.6cm]{\resizebox{#1}{!}{\faIcon{user}}}
\newcommand{\iconglobe}[1][0.6cm]{\resizebox{#1}{!}{\faIcon{globe}}}
\newcommand{\iconclock}[1][0.6cm]{\resizebox{#1}{!}{\faIcon{clock}}}
\newcommand{\iconmail}[1][0.6cm]{\resizebox{#1}{!}{\faIcon{envelope}}}
\newcommand{\icondoc}[1][0.6cm]{\resizebox{#1}{!}{\faIcon{file-alt}}}
\newcommand{\iconsearch}[1][0.6cm]{\resizebox{#1}{!}{\faIcon{search}}}
\newcommand{\icontools}[1][0.6cm]{\resizebox{#1}{!}{\faIcon{tools}}}
\newcommand{\iconchart}[1][0.6cm]{\resizebox{#1}{!}{\faIcon{chart-line}}}
\newcommand{\iconalert}[1][0.6cm]{\resizebox{#1}{!}{\faIcon{bell}}}
\newcommand{\iconlog}[1][0.6cm]{\resizebox{#1}{!}{\faIcon{clipboard-list}}}
\newcommand{\iconeye}[1][0.6cm]{\resizebox{#1}{!}{\faIcon{eye}}}
\newcommand{\iconhacker}[1][0.6cm]{\resizebox{#1}{!}{\faIcon{user-secret}}}
\newcommand{\iconkey}[1][0.6cm]{\resizebox{#1}{!}{\faIcon{key}}}
\newcommand{\iconfile}[1][0.6cm]{\resizebox{#1}{!}{\faIcon{file-alt}}}
\newcommand{\icondog}[1][0.6cm]{\resizebox{#1}{!}{\faIcon{dog}}}
\newcommand{\iconcat}[1][0.6cm]{\resizebox{#1}{!}{\faIcon{cat}}}
\newcommand{\iconvoip}[1][0.6cm]{\resizebox{#1}{!}{\faIcon{phone-volume}}}
\newcommand{\iconmoney}[1][0.6cm]{\resizebox{#1}{!}{\faIcon{money-bill-wave}}}
\newcommand{\icontrash}[1][0.6cm]{\resizebox{#1}{!}{\faIcon{trash}}}
\newcommand{\iconbuilding}[1][0.6cm]{\resizebox{#1}{!}{\faIcon{building}}}

% =============================================================================
% CUSTOM COMMANDS AND ENVIRONMENTS
% =============================================================================

% Key term formatting
\newcommand{\keyterm}[1]{\textbf{#1}}

% Slide section marker
\newcommand{\sectionslide}[1]{
    \begin{frame}
        \vfill
        \centering
        \begin{beamercolorbox}[sep=8pt,center,shadow=true,rounded=true]{title}
            \usebeamerfont{title}\Huge #1
        \end{beamercolorbox}
        \vfill
    \end{frame}
}

% Case study frame environment
\newenvironment{casestudy}[1]{%
    \begin{alertblock}{$\blacktriangleright$ Case Study: #1}
}{%
    \end{alertblock}
}

% Solution frame environment  
\newenvironment{casesolution}[1]{%
    \begin{exampleblock}{$\checkmark$ Solution: #1}
}{%
    \end{exampleblock}
}

% =============================================================================
% FOOTER WITH SLIDE NUMBERS
% =============================================================================
\setbeamertemplate{footline}{
    \leavevmode%
    \hbox{%
        \begin{beamercolorbox}[wd=.333333\paperwidth,ht=2.25ex,dp=1ex,center]{author in head/foot}%
            \usebeamerfont{author in head/foot}\insertshortauthor
        \end{beamercolorbox}%
        \begin{beamercolorbox}[wd=.333333\paperwidth,ht=2.25ex,dp=1ex,center]{title in head/foot}%
            \usebeamerfont{title in head/foot}\insertshorttitle
        \end{beamercolorbox}%
        \begin{beamercolorbox}[wd=.333333\paperwidth,ht=2.25ex,dp=1ex,right]{date in head/foot}%
            \usebeamerfont{date in head/foot}\insertshortdate{}\hspace*{2em}
            \insertframenumber{} / \inserttotalframenumber\hspace*{2ex}
        \end{beamercolorbox}}%
    \vskip0pt%
}

% =============================================================================
% ARTICLE MODE SUPPORT
% =============================================================================
% The following commands help with beamerarticle mode for HTML conversion

% In article mode, add proper section/subsection support and override environments
\mode<article>{
    % Make sure figure captions work properly
    \setbeamertemplate{caption}[numbered]
    \setbeamertemplate{caption label separator}{: }
    
    % Add some spacing for readability
    \setlength{\parskip}{0.5em}
    \setlength{\parindent}{0pt}
    
    % Article mode versions of case study environments
    \renewenvironment{casestudy}[1]{%
        \vspace{0.3cm}\noindent\textbf{Case Study: #1}\newline\begin{itshape}
    }{%
        \end{itshape}\vspace{0.3cm}
    }
    
    \renewenvironment{casesolution}[1]{%
        \vspace{0.3cm}\noindent\textbf{Solution: #1}\newline\begin{itshape}
    }{%
        \end{itshape}\vspace{0.3cm}
    }
}

% Command to mark content that only appears in article mode
\newcommand{\articleonly}[1]{\mode<article>{#1}}

% Command to mark content that only appears in presentation mode
\newcommand{\presentationonly}[1]{\mode<presentation>{#1}}
